\documentclass[11pt,letterpaper]{article}
\usepackage[utf8]{inputenc}

%----- Configuración del estilo del documento------%
\usepackage{epsfig,graphicx}
\usepackage[left=2cm,right=2cm,top=1.8cm,bottom=2.3cm]{geometry}
\usepackage{fancyhdr}
\usepackage{lastpage}
\usepackage{url}
\usepackage{tikz}
\pagestyle{fancy}
\fancyhf{}
\fancyhead[R]{\ifnum\value{page}>1 \textbf{Estadistica Bayesiana}\fi}
\rfoot{\textit{Página \thepage}}
%Encabezado izquierdo: logo
\lfoot{\includegraphics[height=1.0cm]{Logo_FC.png}}
%------ Paquetes matemáticos básicos --------%
\usepackage{amsmath}
\usepackage{amssymb}
\usepackage{amsthm}

\usepackage[spanish]{babel}
\usepackage{graphicx}
\usepackage{hyperref}

\usepackage{tabularx}
\usepackage{xcolor}
\usepackage[table]{xcolor}
\usepackage{colortbl}
\usepackage{array, multirow, multicol, tabularx}
\usepackage{tcolorbox}
\newtheorem{theorem}{Theorem}[section]
\newtheorem{corollary}{Corollary}[theorem]
\newtheorem{lemma}[theorem]{Lemma}

%------si-------%
\definecolor{B}{HTML}{FFFFFF}
\definecolor{G}{HTML}{5e5e5e}
\definecolor{R2}{HTML}{d53d40}
\definecolor{A2}{HTML}{034190}
\definecolor{V2}{HTML}{7faa50}
\newcommand{\R}{\mathbb{R}}
\newcommand{\C}{\mathcal{C}}
\newcommand{\N}{\mathbb{N}}
\newcommand{\Z}{\mathbb{Z}}
\newcommand{\Q}{\mathbb{Q}}
\renewcommand{\theenumi}{\Roman{enumi}}
\renewcommand{\labelenumi}{{\theenumi}.}

\begin{document}
%---Evironment--%
	\makeatletter
        \renewenvironment{proof}[1][\proofname]{\par
            \pushQED{\qed}%
            \normalfont \topsep6\p@\@plus6\p@\relax
            \trivlist
            \item\relax
            {\itshape
            #1\@addpunct{.}}\par\vspace{\baselineskip}\ignorespaces
            }{%
            \popQED\endtrivlist\@endpefalse
            }
        \newenvironment{Solucion}[1][Soluci\'on]{\par
            \pushQED{\qed}%
            \normalfont \topsep6\p@\@plus6\p@\relax
            \trivlist
            \item\relax
            {\itshape
            #1\@addpunct{.}}\par\vspace{\baselineskip}\ignorespaces
            }{%
            \popQED\endtrivlist\@endpefalse
            }
            
    \makeatother

%------ Encabezado -------- %

\begin{center}
    \begin{minipage}{3cm}
    	\begin{center}
    		\includegraphics[height=3.4cm]{logo_unam.png}
    	\end{center}
    \end{minipage}\hfill
    \begin{minipage}{10cm}
    	\begin{center}
    	\textbf{\large Universidad Nacional Autónoma de México}\\[0.1cm]
        \textbf{Facultad de Ciencias}\\[0.1cm]
        \textbf{\'Estadistica Bayesiana}\\[0.1cm]
        Tareita 6 \\[0.1cm]
         El\'ias L\'opez Rivera\\[0.1cm]
        \texttt{ elias.lopezr,@ciencias.unam.mx }\\[0.1cm]
        Fecha:\,\,24/09/2026
    	\end{center}
    \end{minipage}\hfill
    \begin{minipage}{3cm}
    	\begin{center}
    		\includegraphics[height=3.4cm]{Logo_FC.png}
    	\end{center}
    \end{minipage}
\end{center}
\rule{17cm}{0.1mm}
%------ Fin de encabezado -------- %
\begin{tcolorbox}[
	title = \textcolor{black}{\textcolor{white}{Problema 1}},]
\textit{Dado el valor esperado $\mu=0.4$ y varianza $\sigma=0.02$, obtener las distribuciones iniciales del parámetro y predictiva de un modelo Beta-Binomial. Considerando un tamaño de muestra m=18 (creo que era 18), y utilizando solamente las primeras 13 observaciones, obtener las distribuciones finales del parámetro y predictiva, sus intervalos de confianza, los posibles estimadores del parámetro (como medidas de tendencia central), observar si las últimas 5 observaciones se encuentran en estos intervalos de confianza, y hacer cualquier otra observación o experimento sobre el ejercicio.
}
\end{tcolorbox}
\textbf{Soluci\'on}
Primero recordemos que suponemos $\theta\sim\beta(\alpha,\beta)$ con $\mu=\frac{\alpha}{\alpha+\beta}$ y $\sigma^2=\frac{\alpha\beta}{(\alpha+\beta)^2(\alpha+\beta+1)}$, tomamos $f(x_i|\theta)\sim bin(n=18,\theta)$, para $x_i\in\{1,2,\cdots,18\}$ y solo
tomamor las primeras $13$ observaciones $S_{13}$\,\\
\,\\
Recordemos que la posteriori de $\theta$ dado $s_{13}$ se distrbuye beta con parametros $\alpha'=\alpha+Y$, donde $Y$ es el n\'umero de exitos, $\beta'=\beta+\bar{Y}$ donde $\bar{Y}$ es el numero de fracasos:\,\\
\begin{equation*}
    f(\theta|S_{13})=\frac{\theta^{\alpha'-1}(1-\theta)^{\beta'-1}}{\beta(\alpha',\beta')}
\end{equation*}
Finalmente la distribucion predictiva se obtiene como una beta binomial:\,\\
\begin{equation*}
    P(x|s_{13})=\binom{18}{x}\frac{\beta(x+\alpha',18-x+\beta')}{\beta(\alpha',\beta')}
\end{equation*}
Simulamos las 18 observaciones obtenemos el vector de resultados:\,\,\,[ 6  6  5 11  9 10  4 12 10  7  6  6  7  8  8  7  9  6]
\begin{center}
    \includegraphics[height=0.4\textwidth]{Captura de pantalla (830).png}
\end{center}
Con esto en mente calculamos los datos de cada distribuci\'on:\,\\
\begin{equation*}
    \beta=\alpha\left(\frac{1-0.4}{0.4}\right)=1.5\alpha
\end{equation*}
\begin{align*}
    0.02&=\frac{1.5\alpha^2}{(2.5)^2\alpha^2(2.5\alpha+1)}=\frac{1.5}{(6.25)(2.5\alpha+1)}\,\\
    \,\\
    &\implies 0.02(6.25)(2.5\alpha+1)=1.5\implies 2.5\alpha=11\implies \alpha=4.4\implies \beta=6.6
\end{align*}
Por tanto $\theta\sim\beta(4.4,6.6)$, ahora encontramos $Y$ e $\bar{Y}$:\,\\
\begin{equation*}
    Y=6+6+5+11+9+10+4+12++10+7+6+6+7=99\implies \alpha'=4.4+99=103.4
\end{equation*}
\begin{equation*}
    \bar{Y}=13*18-99=135\implies \beta'=6.6+135=141.6
\end{equation*}
Por tanto la posteriori nos da $\theta|s_13\sim\beta(103.4,141.6)$, calculamos los estimadores $E[\theta|s_{13}],\theta_{mediana},\theta_{moda}$, como la distribuci\'on es una
beta tenemos formulas conocidad para ellas:\,\\
\begin{align*}
    \theta_{media}&=\frac{\alpha'}{\alpha'+\beta'}=\frac{103.4}{103.4+141.6}\approx0.4220\,\\
    \,\\
    \theta_{moda}&=\frac{\alpha'-1}{\alpha'+\beta'-2}=\frac{103.4-1}{103.4+141.6-2}\approx0.4214
\end{align*}
Calulamos el cuantil de $\frac{1}{2}$ para la mediana $\beta(F^{-1}(0.5))=0.418$\,\\
\,\\
Finalmente calculamos un intervalo de confianza al $0.85$, dividimos el error $(0.15)$ simetricamente por tanto encontramos los cuantiles
$I_1=F^{-1}(\frac{0.15}{2}=0.075)=0.3773$ y $I_2=F^{-1}(1-\frac{0.15}{2}=0.925)=0.4667$, por tanto nuestro intervalo de confianza es $[I_1,I_2]$ claramente el valor original
simulado cumple $\theta_0\in[I_1,I_2]$, ademas que los otros estimadores claramente se acercaron bastante al valor original\,\\
\,\\
Ahora calculamos los valores de la predictiva posteriori con los valores de la muestra restantes $[8,8,7,9,6]$:\,\\
\begin{equation*}
    p(x=8|s_{13})=\binom{18}{8}\frac{\beta(18+103.4,18-8+141.6)}{\beta(103.4,141.6)}\approx0.1200
\end{equation*}
\begin{equation*}
    p(x=7|s_{13})=\binom{18}{7}\frac{\beta(18+103.4,18-7+141.6)}{\beta(103.4,141.6)}\approx0.01415
\end{equation*}\,\\
\begin{equation*}
    p(x=9|s_{13})=\binom{18}{9}\frac{\beta(18+103.4,18-9+141.6)}{\beta(103.4,141.6)}\approx0.0898
\end{equation*}\,\\
\begin{equation*}
    p(x=6|s_{13})=\binom{18}{6}\frac{\beta(18+103.4,18-6+141.6)}{\beta(103.4,141.6)}\approx0.1472
\end{equation*}\,\\
Finalmente calculamos un intervalo de confianza de $0.85$, de nuevo dividimos el error $(0.15)$ simetricamente, por tanto encontramos los cuantiles
$J_1=F^{-1}_{post}(0.075)=5$, $J_2=F^{-1}_{post}(0.925)=11$, por tanto nuestro inervalo de confianza es $[5,11]$, podemos notar que clamente todos los valores
posteriores de la muestra est\'an dentro $\{8,8,7,9,6\}\in[5,11]$, concluimos que el metodo funciona de manera m\'as que eficiente
\end{document}